void loop()
{
    // 测姿态-----------------------------------------------------------------------------------------------------------------------------------------------
    // MPU6050_check();

    // 测距离------------------------------------------------------------------------------------------------------------------------------------------------------
    // float distance = Distance_chect();
    // Serial.print("距离：");
    // Serial.print(distance);
    // Serial.println("cm");
    // float distance = 0;
    // netPrintf("距离：%fcm\n", distance);

    // 测速度-----------------------------------------------------------------------------------------------------------------------------------------------
    // Encoder_Check();
    Velocity_Check();

    // 测电机---------------------------------------------------------------------------------------------------------------------------------
    //  my_motor[0].Motor_Run(1,50);
    //  my_motor[1].Motor_Run(2,50);
    //  my_motor[2].Motor_Run(3,50);
    //  my_motor[3].Motor_Run(4,50);

    //  delay(3000);

    //  my_motor[0].Motor_Run(1,0);
    //  my_motor[1].Motor_Run(2,0);
    //  my_motor[2].Motor_Run(3,0);
    //  my_motor[3].Motor_Run(4,0);

    //  delay(3000);

    // Motor_test1(100);

    // 测PID-------------------------------------------------------------------------------------------------------------------------------
    // Pid_controller_run(400.0);//单位mm/s

    // for (int i = 0; i < 4; i++)
    // {
    //     run[i].Pid_run_();
    // }

    // 测运动学解算---------------------------------------------------------------------------------------------------------------------------
    // float target_linear = 100;
    // float target_angular = 0.5;
    // float out_MotorL, out_MotorR;
    // Kinematics_down(target_linear, target_angular, &out_MotorL, &out_MotorR);
    // run[0].Pid_init(1, out_MotorL);
    // run[1].Pid_init(2, out_MotorL);
    // run[2].Pid_init(3, out_MotorR);
    // run[3].Pid_init(4, out_MotorR);

    run[0].Pid_run_();
    run[1].Pid_run_();
    run[2].Pid_run_();
    run[3].Pid_run_();

    // 查看实际给车子的脉冲的占空比
    //  a=pid_controller[1].update(Velocity[0]);
    //  b=pid_controller[2].update(Velocity[1]);
    //  c=pid_controller[3].update(Velocity[2]);
    //  d=pid_controller[4].update(Velocity[3]);
    //  Serial.printf("给电机的pwm：%f,%f,%f,%f\n",a,b,c,d);

    // 里程计打印--------------------------------------------------------------------------------------------------------------------------------
    odom_update();

}
